\documentclass[11pt,a4paper]{article}
\usepackage[utf8]{inputenc}
\usepackage[T1]{fontenc}
\usepackage{lmodern}
\usepackage[french]{babel}
\usepackage{amsmath,amssymb,mathtools}
\usepackage{icomma}
\usepackage{xcolor}
\usepackage[most]{tcolorbox}
\usepackage{enumitem}
\usepackage[a4paper,margin=2.2cm,headheight=26pt,headsep=10pt]{geometry}
\usepackage{fancyhdr}
\usepackage[hidelinks]{hyperref}
\definecolor{primary}{RGB}{222,49,99}
\definecolor{primarydark}{RGB}{150,22,55}
\definecolor{excol}{RGB}{0,135,90}
\tcbset{boxbase/.style={enhanced,breakable,boxrule=0.9pt,arc=2.5pt,
  left=3mm,right=3mm,top=2.3mm,bottom=2.3mm,fonttitle=\bfseries,coltitle=white}}
\newtcolorbox[auto counter]{corrige}[1][]{boxbase,colback=excol!4,colframe=excol,
  title={\textsc{Corrigé}~\thetcbcounter\ifx\relax#1\relax\relax\else\textmd{ --- \itshape #1}\fi}}
\newcommand{\R}{\mathbb{R}}
\newcommand{\N}{\mathbb{N}}
\newcommand{\ds}{\displaystyle}
\pagestyle{fancy}\fancyhf{}
\fancyhead[L]{\small\textcolor{primarydark}{\textbf{Thème 1} — Puissances, racines et exposants}}
\fancyhead[R]{\small\textcolor{primarydark}{\textsc{Corrigé — Difficile}}}
\fancyfoot[C]{\small\thepage}
\renewcommand{\headrulewidth}{0.8pt}
\begin{document}

\begin{center}
{\LARGE\bfseries\color{primarydark} Niveau Difficile — Corrigé}\\[2pt]
{\color{primary}\rule{0.45\textwidth}{1.2pt}}
\end{center}
\vspace{1mm}

\begin{corrige}[exposants fractionnaires]
\textbf{a)} On convertit chaque racine en exposant fractionnaire, puis on multiplie les exposants emboîtés :
\[ \sqrt[n]{\sqrt[n]{c^{2n^{2}}}}=\Bigl(\bigl(c^{2n^{2}}\bigr)^{1/n}\Bigr)^{1/n}
   =\bigl(c^{2n^{2}}\bigr)^{1/n^{2}}=c^{\,2n^{2}\cdot\frac{1}{n^{2}}}=c^{2}. \]
Le résultat ne dépend pas de $n$ car l'exposant $\dfrac{2n^{2}}{n^{2}}=2$ se simplifie entièrement.

\smallskip
\textbf{b)} $\ds E=\frac{c^{1/3}\,c^{5/6}}{c^{1/2}}=c^{\frac13+\frac56-\frac12}
=c^{\frac{2}{6}+\frac{5}{6}-\frac{3}{6}}=c^{\frac46}=c^{2/3}.$

\smallskip
\textbf{c)} Pour $c=64$ : $E=64^{2/3}=\bigl(64^{1/3}\bigr)^{2}=4^{2}=16$ (car $64^{1/3}=\sqrt[3]{64}=4$).
\end{corrige}

\begin{corrige}[équations exponentielles]
\begin{align*}
\text{a) }&2x+1=-x\Rightarrow 3x=-1\Rightarrow x=-\tfrac13, &&S=\bigl\{-\tfrac13\bigr\};\\
\text{b) }&4^{x+1}=2^{2x+2},\ 8^{x}=2^{3x}\Rightarrow 2x+2=3x\Rightarrow x=2, &&S=\{2\};\\
\text{c) }&\frac{2^{x}2^{x+1}}{2^{3}}=2^{2x+1-3}=2^{2x-2}=2^{0}\Rightarrow 2x-2=0\Rightarrow x=1, &&S=\{1\};\\
\text{d) }&\Bigl(\tfrac{9^{x}}{3}\Bigr)^{2}=\bigl(3^{2x-1}\bigr)^{2}=3^{4x-2}=3^{6x-4}\Rightarrow 4x-2=6x-4\Rightarrow x=1, &&S=\{1\}.
\end{align*}
\end{corrige}

\begin{corrige}[dilutions successives]
\textbf{a)} $\ds C(3)=C_0\Bigl(\frac12\Bigr)^{3}=\frac{C_0}{8}.$

\smallskip
\textbf{b)} On résout $\ds C_0\Bigl(\frac12\Bigr)^{x}=\frac{C_0}{32}$, soit (en divisant par $C_0>0$)
$\ds\Bigl(\frac12\Bigr)^{x}=\frac{1}{32}=\Bigl(\frac12\Bigr)^{5}\Rightarrow x=5.$ Il faut \textbf{5 étapes}.

\smallskip
\textbf{c)} $\ds C(x)\leq\frac{C_0}{64}\Leftrightarrow\Bigl(\frac12\Bigr)^{x}\leq\Bigl(\frac12\Bigr)^{6}$.
La base $\tfrac12<1$ rend $x\mapsto\bigl(\tfrac12\bigr)^{x}$ \emph{décroissante} : l'inégalité sur les
valeurs \emph{s'inverse} sur les exposants, donc $x\geq 6$. La condition est atteinte \textbf{à partir de la $6^{\text{e}}$ étape}.
\end{corrige}

\end{document}
